\begin{abstract}
Over the past 7 years, attention has become one of the most important primitives in deep learning. The primary approach to optimize attention is FlashAttention, which fuses the operation together, drastically improving both the runtime and the memory consumption. However, the importance of FlashAttention combined with its monolithic nature poses a problem for researchers aiming to try new attention variants --- a ``software lottery". 
This problem is exacerbated by the difficulty of writing efficient fused attention kernels, resisting traditional compiler-based approaches. 
We introduce \textit{FlexAttention}, a novel compiler-driven programming model that allows implementing the majority of attention variants in a few lines of idiomatic PyTorch code. We demonstrate that many existing attention variants (e.g. Alibi, Document Masking, PagedAttention, etc.) can be implemented via \textit{FlexAttention}, and that we achieve competitive performance compared to these handwritten kernels. Finally, we demonstrate how \textit{FlexAttention} allows for easy \textit{composition} of attention variants, solving the combinatorial explosion of attention variants.
\end{abstract}